\section{Harmonic measures and the finite-zero calculation}
\label{sec:harmonic}
This section proves the zero-side statements independently of an
arithmetic cutoff. Put $\mathcal D= C_c^\infty(\R)$ and let
$\mathcal S(\R)$ be the Schwartz space, defined by the seminorms
$\sup_t(1+|t|)^j|f^{(k)}(t)|$ for all nonnegative integers $j,k$.
We use logarithmic Haar measure $dt=dx/x$ on $\R_+^\times$,
convolution $(f*g)(t)=\int f(u)g(t-u)du$, and
$f^*(t)=\overline{f(-t)}$. The transform
$\widehat f(\tau)=\int f(t)e^{i\tau t}dt$ is unitary from
$L^2(\R,dt)$ to $L^2(\R,d\tau/(2\pi))$.
For a bounded measurable $a$, let $C_a$ denote its Fourier multiplier.
For $h\in\mathcal D$ put $H(z)=\int h(t)e^{zt}dt$.

A zero divisor here is a locally finite multiset
$\mathcal Z=\{(a_\rho,\gamma_\rho,m_\rho)\}$, where
$|a_\rho|\leq d$ for a fixed $d$, the ordinates are real, and the
positive integer $m_\rho$ is a multiplicity. We assume
\begin{equation}\label{eq:weighted-count}
 Z_2:=\sum_{\rho\in\mathcal Z}\frac{m_\rho}{1+\gamma_\rho^2}<\infty.
\end{equation}
The usual $O(T\log(T+2))$ zero count for a fixed primitive Dirichlet
or Hecke $L$-function implies this condition by partial summation.
Define
\[
 d\mu_{a,\gamma}(\tau)=
 \begin{cases}
 \displaystyle\frac{|a|}{\pi((\tau-\gamma)^2+a^2)}\,d\tau,&a\ne0,\\
 \delta_\gamma(d\tau),&a=0.
 \end{cases}
\]
Its Fourier transform is $e^{-|a||t|}e^{i\gamma t}$.

\begin{proposition}\label{prop:summability}
The measure $\mathfrak b=\sum_\rho m_\rho\mu_{a_\rho,\gamma_\rho}$
is locally finite and integrates the absolute value of every Schwartz
function. Both $\mathfrak B(h)=\int H(i\tau)d\mathfrak b(\tau)$ and
$\mathfrak W(h)=\sum_\rho m_\rho H(a_\rho+i\gamma_\rho)$ converge
absolutely for $h\in\mathcal D$. Moreover
$\mathfrak B(g*g^*)\geq0$ for $g\in\mathcal D$.
\end{proposition}
\begin{proof}
For $\phi\in\mathcal S$, uniformly in $0<b\leq d$,
\[
 \int |\phi(\tau)|\frac{b\,d\tau}{\pi((\tau-\gamma)^2+b^2)}
 \leq \frac{C_{d,\phi}}{1+\gamma^2}.
\]
For $|\gamma|\geq2$, split the integral into
$|\tau-\gamma|\leq|\gamma|/2$ and its complement. In the first
part, rapid decrease gives $|\phi(\tau)|\leq C/\gamma^2$ and the
Poisson mass is at most one. In the second, the density is at most
$4d/(\pi\gamma^2)$ and $\phi$ is integrable. For bounded $\gamma$
use its supremum. The same estimate holds for a Dirac mass.
It proves the claimed summability by \eqref{eq:weighted-count} and,
using a Schwartz function bounded below on a chosen compact interval,
also proves local finiteness. Integration by parts twice in $H$
gives a bound $C_h/(1+\gamma^2)$ uniformly for $|a|\leq d$.
Finally $\widehat{g*g^*}(\tau)=|\widehat g(\tau)|^2$.
\end{proof}

\begin{theorem}[The finite-packet formula]\label{thm:finite-packet}
Let $F$ be a finite set of distinct numbers $z=a+i\gamma$, with
assigned multiplicities $m_z$.
In $L^2(\R)$ let $V_L$ be the span of
\[
 1_{[-L,L]}(t)t^j e^{\overline zt},\qquad z\in F,\quad 0\leq j<m_z,
\]
and let $P_{F,L}$ be its orthogonal projection. For every bounded
continuous $\phi:\R\to\C$,
\begin{equation}\label{eq:packet-limit}
 \lim_{L\to\infty}\Tr(P_{F,L}C_\phi)
       =\sum_{z\in F}m_z\int_\R\phi(\tau)d\mu_{a,\gamma}(\tau).
\end{equation}
There is no zero-location or simplicity assumption.
\end{theorem}
\begin{proof}
Replacing $t^j$ by $(t-L)^j$ or $(t+L)^j$ does not change the
span at a fixed $z$. For $a>0$, translate the upper endpoint to
zero and remove the scalar exponential and phase. The resulting
vectors are
$1_{[-2L,0]}(u)u^j e^{au-i\gamma u}$, which converge in $L^2$
to $1_{(-\infty,0]}(u)u^j e^{au-i\gamma u}$.
For $a<0$, the analogous lower-endpoint vectors converge to
$1_{[0,\infty)}(u)u^j e^{-|a|u-i\gamma u}$.
Each finite limiting family is linearly independent: an exponential
polynomial that vanishes on an interval vanishes identically, and
successive application of the differential factors annihilating all
but one exponent proves each coefficient zero. Its Gram matrix is
therefore positive definite.

For $a=0$ use the scaled vectors
$(2L)^{-1/2}1_{[-L,L]}(t)((t+L)/(2L))^j e^{-i\gamma t}$.
For one ordinate their Gram matrix is
$(1/(j+k+1))_{j,k}$, which is positive definite because it is the
Gram matrix of monomials in $L^2(0,1)$. Between distinct ordinates
the entries tend to zero by integration by parts after scaling.
The scalar Fourier measures of these vectors concentrate at
$\gamma$: scaling writes the transformed profile as a fixed
$L^2$ function evaluated at $2L(\tau-\gamma)$.
Consequently $(C_\phi-\phi(\gamma))v_{j,L}\to0$ in $L^2$,
first by continuity near $\gamma$ and then by the small Fourier
tail of the fixed profile. Thus this block contributes
$m_z\phi(\gamma)$.

The inner products of an endpoint vector with a line vector are
$O(L^{-1/2})$, by the exponential integrability of every endpoint
polynomial. Opposite endpoint families have vanishing inner
products, by translating one of two fixed $L^2$ functions a
distance $2L$ and approximating them by compactly supported functions.
Therefore the Gram matrix of all normalized vectors tends to the
block diagonal matrix of the two endpoint families and the separate
line families. Its inverse square root converges. Orthonormalizing
with this inverse square root changes the trace by $o(1)$ relative
to the sum of the block traces, since the dimension is fixed and
$\|C_\phi\|\leq\|\phi\|_\infty$.

It remains to compute one endpoint trace. For the upper endpoint
the transforms span the rational functions
$(a+i(\tau-\gamma))^{-j-1}$. Set $\lambda=a+i\gamma$ in the
right half-plane and
\[
 b_\lambda(w)=\frac{w-\lambda}{w+\overline\lambda},\qquad
 k_\lambda(w)=\frac{\sqrt{2a}}{w+\overline\lambda}.
\]
List the zeros with repetitions as $\lambda_1,\ldots,\lambda_n$.
An orthonormal basis is
\[
 e_j(w)=k_{\lambda_j}(w)\prod_{r<j}b_{\lambda_r}(w).
\]
Here is the finite-dimensional justification. The kernel
$k_\lambda$ has norm one by direct integration on $w=i\tau$.
It reproduces evaluation at $\lambda$, up to the factor
$\sqrt{2a}$, by Cauchy's integral formula. Multiplication by
$b_\lambda$ is an isometry on boundary values. Its range consists
of the analytic functions vanishing at $\lambda$: division removes
the zero and preserves the boundary norm. Induction gives the
orthogonal decomposition by the displayed $e_j$. Partial fractions
show that its finite-dimensional complement is exactly the span of
the stated poles, including the repeated orders. This also proves
the asserted basis without assuming distinct ordinates.

Since $|b_\lambda(i\tau)|=1$, its diagonal density is
\[
 \sum_{j=1}^n |e_j(i\tau)|^2\frac{d\tau}{2\pi}
 =\sum_{j=1}^n\frac{a_j\,d\tau}
               {\pi(a_j^2+(\tau-\gamma_j)^2)}.
\]
The trace against $C_\phi$ is its integral against $\phi$.
Reflection proves the same formula for the lower endpoint, with
$a$ replaced by $|a|$. Combining the blocks proves
\eqref{eq:packet-limit}.
\end{proof}

\begin{remark}
For two simple upper-endpoint zeros with the same ordinate and
positive real parts $a,b$, their individually normalized vectors
have limiting inner product $2\sqrt{ab}/(a+b)$. It is nonzero.
The finite Blaschke computation is what permits the trace to be
additive despite this nonorthogonality. This supplies the needed
Gram calculation in the auxiliary computation following Connes's
Lemma 3, where asymptotic orthogonality is stated too broadly.
\end{remark}

\section{The discrepancy and a zero-free half-plane}
This section states exactly what zero location controls on the zero
side of the formula. It does not assume an identification with a
prolate trace.

\begin{proposition}[There is no per-pair sign]\label{prop:sign}
For $a>0$ and $\gamma\in\R$, let $\mathfrak B_{a,\gamma}$ and
$\mathfrak W_{a,\gamma}$ be the contributions of the reflected
pair $a+i\gamma,-a+i\gamma$, each with multiplicity one.
Then
\begin{equation}\label{eq:pair-discrepancy}
 (\mathfrak B_{a,\gamma}-\mathfrak W_{a,\gamma})(h)
       =-2\int_\R h(t)\sinh(a|t|)e^{i\gamma t}\,dt.
\end{equation}
This functional takes both positive and negative values on tests
$h=g*g^*$ with $g\in\mathcal D$.
\end{proposition}
\begin{proof}
The Poisson identity gives the pair kernel $2e^{-a|t|}e^{i\gamma t}$.
The two Weil evaluations give $2\cosh(at)e^{i\gamma t}$.
Subtracting proves \eqref{eq:pair-discrepancy}.
Choose a real nonnegative smooth bump $\eta$ of integral one and
compact support, put $\eta_\varepsilon(t)=\varepsilon^{-1}
\eta(t/\varepsilon)$, and fix $d>0$. For
\[
 g_\varepsilon^\pm(t)=e^{-i\gamma t}
       (\eta_\varepsilon(t)\mathbin\pm\eta_\varepsilon(t-d)),
\]
the double-integral form of \eqref{eq:pair-discrepancy} and
continuity of $\sinh(a|t-s|)$ show that its values tend respectively
to $-4\sinh(ad)$ and $4\sinh(ad)$. The two diagonal bump terms
tend to zero and the two cross terms give these limits. For small
positive $\varepsilon$ the signs are therefore opposite.
\end{proof}

\begin{theorem}[Quantitative band estimate]\label{thm:band-bound}
Assume \eqref{eq:weighted-count} and $|a_\rho|\leq\delta$.
For $h\in\mathcal D$ define
\[
 \mathcal N_\delta(h)=|h(0)|+
  \int_\R e^{\delta|t|}
  \bigl(|t||h''(t)|+2|h'(t)|+(1+\delta^2)|t||h(t)|\bigr)dt.
\]
Then
\begin{equation}\label{eq:band-bound}
 |\mathfrak B(h)-\mathfrak W(h)|
 \leq 2\mathcal N_\delta(h)
          \sum_\rho\frac{m_\rho|a_\rho|}{1+\gamma_\rho^2}
 \leq2\delta\mathcal N_\delta(h)Z_2.
\end{equation}
The estimate assumes the stated band and no RH. For the usual
functional-equation-symmetric zero divisor, a zero-free half-plane
$\Re s>\theta$ gives $\delta=\theta-1/2$.
\end{theorem}
\begin{proof}
For one zero put
$e_a(t)=h(t)(e^{-|a||t|}-e^{at})$. For $a>0$ this is
$-2h(t)\sinh(at)1_{t>0}$; for $a<0$ it is
$2h(t)\sinh(|a|t)1_{t<0}$. It is continuous at zero. Its
first derivative has jump $-2|a|h(0)$, so its distributional
second derivative contains $-2|a|h(0)\delta_0$.
On the open half-line where it is nonzero, differentiation twice
and the inequalities
$\sinh(|a||t|)\leq|a||t|e^{\delta|t|}$ and
$\cosh(|a|t)\leq e^{\delta|t|}$ give
\[
 \|e_a\|_1+\|e_a''\|_{\mathrm{TV}}
      \leq2|a|\mathcal N_\delta(h).
\]
Because the Fourier transform of $e_a-e_a''$ is
$(1+\gamma^2)\widehat e_a(\gamma)$, this proves
$|\widehat e_a(\gamma)|\leq
2|a|\mathcal N_\delta(h)/(1+\gamma^2)$.
Summing the absolutely convergent estimates proves the theorem.
The functional equation reflects $a$ to $-a$ with the same
ordinate in a fixed unitary sector (or, equivalently, after
combining conjugation with the dual-character equation).
\end{proof}
